\begin{claim}{C.84}{Proved}{The exact radial translation topology and a contractible monotone subspace}
Retain the actual radial coordinate
$j(z)=z/(1+\|z\|)$, its inverse, the original
closed unit disk, and the operator
$q=jT_{(1,0)}j^{-1}$ of C.83.
Conjugation $h\mapsto j^{-1}hj$ is a
topological group isomorphism from the
actual relative centralizer $\mathscr C_Q$
onto the following explicitly topologized
group $\mathscr P$.
Its elements are plane homeomorphisms
$F=(X_F,Y_F)$ commuting with $T_{(1,0)}$
and satisfying
\[
 \lim_{R\to\infty}
  \sup_{\substack{0\leq x\leq1\\ |y|\geq R}}
   \frac{\|F(x,y)-(x,y)\|}{|y|}=0.
                                                               \tag{C.84.1}
\]
Every inverse satisfies the same condition.
The exact topology on $\mathscr P$ is
given by the finite metric
\[
 \begin{gathered}
 \|F-H\|_*=
  \sup_{\substack{0\leq x\leq1\\ y\in\mathbb R}}
     \frac{\|F(x,y)-H(x,y)\|}{1+|y|},\\
 \delta_*(F,H)=\|F-H\|_*+\|F^{-1}-H^{-1}\|_*.
 \end{gathered}                                                  \tag{C.84.2}
\]
In (C.84.1) the limit is uniform over
every compact family in the actual group.

Let $\mathscr M\subset\mathscr P$ consist
of those $F$ for which, at every $x$,
$y\mapsto Y_F(x,y)$ is an increasing
homeomorphism of $\mathbb R$.
The proof gives an explicit continuous
contraction of this entire subspace,
fixing the identity, through the same
subspace. Conjugating back through $j$
gives the contraction in the original
closed-disk topology, preserving its
circle and every original operator label.
The inclusion $\mathscr M\hookrightarrow\mathscr P$
is proper, as witnessed by an explicit
periodic core rotation below.
\end{claim}

\begin{proof}
\emph{The exact set correspondence and its compact-family bound.}
For a commuting plane homeomorphism write
\[
 F(x,y)=(x+u_F(x,y),\,y+w_F(x,y)).
                                                               \tag{C.84.3}
\]
Both displacement functions are periodic
of period one in $x$. We first prove
necessity of (C.84.1) uniformly over a
compact family of actual disk maps.

Take any sequence $F_i$ from such a
family and points $(x_i,y_i)$ with
$0\leq x_i\leq1$ and $|y_i|\to\infty$.
Choose integers $m_i^+,m_i^-$ such that
\[
 x_i+m_i^+=y_i+e_i^+,\qquad
 x_i+m_i^-=-y_i+e_i^-,\qquad
 |e_i^\pm|\leq\tfrac12.
                                                               \tag{C.84.4}
\]
Pass to a subsequence on which the sign
of $y_i$ is constant. The two radial
input sequences approach the corresponding
two diagonal points of the original
circle. Compactness of the actual family,
whose members fix every circle point,
forces the output sequences to approach
those same points. To justify the uniform
assertion directly, any alleged failure
has a subsequence of disk maps converging
uniformly; its limiting map fixes the
specified boundary point.

Periodicity makes the two output second
coordinates equal to $Y_i=Y_{F_i}(x_i,y_i)$.
The limiting directions have nonzero
second coordinate. Thus the following
ratios are defined eventually and satisfy
\[
 \frac{X_{F_i}(x_i,y_i)+m_i^+}{Y_i}\longrightarrow1,
 \qquad
 \frac{X_{F_i}(x_i,y_i)+m_i^-}{Y_i}\longrightarrow-1.
                                                               \tag{C.84.5}
\]
Subtracting and using
$(m_i^+-m_i^-)/y_i\to2$ gives $Y_i/y_i\to1$.
Either ratio then gives
$u_{F_i}(x_i,y_i)/y_i\to0$.
Consequently the normalized norm in
(C.84.1) tends to zero along every such
sequence. This proves the required
uniform compact-family limit.
Apply the same proof to the compact
family of actual inverse maps.

Conversely, suppose a commuting plane
homeomorphism satisfies (C.84.1).
For any $\epsilon>0$ its displacement
is at most $\epsilon|y|$ outside one
horizontal band, uniformly in all $x$
by periodicity. Within that band the
displacement has a finite bound, by
continuity on its compact quotient.
It follows that
\[
 \frac{\|F(z)-z\|}{1+\|z\|}\longrightarrow0
                   \quad(\|z\|\to\infty).
                                                               \tag{C.84.6}
\]
The exact radial estimate (C.83.3)
therefore proves that $jFj^{-1}$
extends continuously by the identity
to the original circle. This extension
is a continuous bijection of the
compact disk onto itself, so its
inverse is continuous. It is an
actual member of $\mathscr C_Q$.
Applying necessity to that inverse
also proves its version of (C.84.1).
Restriction and extension by the
specified identity are mutually inverse,
and conjugation preserves multiplication
and every original power of $q$.
This proves the set and group
correspondence without assuming a
different compactification.

\emph{Both inverse topology maps.}
The expressions in (C.84.2) are finite:
on compact $y$-bands they are bounded
by periodicity, and on the tails this
follows from (C.84.1).
The displayed formula is a metric.

Suppose actual disk maps converge,
with their inverses, to an actual map.
The sequence together with its limit
is a compact family, so the uniform
version of (C.84.1) bounds both weighted
tails in (C.84.2). On each compact
rectangle $[0,1]\times[-R,R]$,
conjugation through $j$ and $j^{-1}$
gives uniform convergence in plane
coordinates. The images remain in
one compact interior set, where
$j^{-1}$ is uniformly continuous.
Combining this compact-rectangle
convergence with the tail bound proves
convergence of $\|F_i-F\|_*$.
The same proof for inverses gives
convergence in $\delta_*$.

For the other direction fix
$H\in\mathscr P$ and set
\[
 c_H=\sup_{(x,y)\in\mathbb R^2}
       \frac{1+|y|}{1+\|H(x,y)\|}<\infty.
                                                               \tag{C.84.7}
\]
The finiteness follows from
$|Y_H(x,y)|\geq|y|/2$ on sufficiently
far tails, and boundedness of the
numerator on the remaining band.
Apply (C.83.3) at the base point
$H(x,y)$ to obtain
\[
 \sup_{z\in\mathbb R^2}\|jF(z)-jH(z)\|
             \leq2c_H\|F-H\|_*.
                                                               \tag{C.84.8}
\]
The periodicity of $F-H$ permits
exactly the supremum in (C.84.2).
The inverse version of (C.84.8)
proves convergence in the original
disk map-and-inverse metric.
These two metrizable-space arguments
prove the asserted topological group
isomorphism.

\emph{Exact factorization on the monotone subspace.}
For $F\in\mathscr M$ define
\[
 \begin{gathered}
 V_F(x,y)=(x,Y_F(x,y)),\qquad
 H_F=F V_F^{-1},\\
 H_F(x,y)=(h_{F,y}(x),y),\qquad F=H_F V_F.
 \end{gathered}                                                  \tag{C.84.9}
\]
The inverse of $V_F$ is continuous.
Indeed, for a converging sequence of
$x$-coordinates and target heights,
two fixed height values bracketing
the limiting inverse bracket all
nearby inverses. Every convergent
subsequence between those brackets
must solve the same strictly monotone
equation. Shrinking the brackets
proves convergence of the inverse.
This gives an actual plane homeomorphism.
Its commutation with $T_{(1,0)}$
is literal periodicity, and its
displacement satisfies (C.84.1).
Thus $V_F\in\mathscr P$, and so
$H_F\in\mathscr P$ as well.

The formula for $H_F$ fixes the height.
For each height its first coordinate
is a real-line homeomorphism and satisfies
$h_{F,y}(x+1)=h_{F,y}(x)+1$.
It is therefore increasing.
For $0\leq r\leq1$ put
\[
 \begin{split}
 V_{F,r}(x,y)
    &=(x,(1-r)y+rY_F(x,y)),\\
 H_{F,r}(x,y)
    &=((1-r)x+rh_{F,y}(x),y),\\
 A_r(F)&=H_{F,r}V_{F,r}.
 \end{split}                                                    \tag{C.84.10}
\]
Every displayed one-dimensional
convex combination is strictly
increasing and proper, and its
inverse depends continuously on
the other coordinates by the
same bracketing argument.
The maps retain the translation
equations exactly. Their
displacements are scalar multiples
of the respective original fiber
displacements, so (C.84.1) holds
uniformly for $r\in[0,1]$.
They are actual elements of
$\mathscr P$.

The second coordinate of $A_r(F)$
is the increasing function
$(1-r)y+rY_F(x,y)$.
Thus $A_r(F)\in\mathscr M$.
The faces satisfy
$A_0(F)=1$, $A_1(F)=F$, and
$A_r(1)=1$.

\emph{Full family continuity, including inverses.}
The map $F\mapsto V_F$ is continuous
in weighted forward distance, by
its formula and (C.84.2).
Equation (C.84.8) converts this
to uniform convergence of the
actual disk maps.
For homeomorphisms of this compact
disk, uniform forward convergence
to a homeomorphism implies uniform
inverse convergence. Explicitly,
if $\omega_{B^{-1}}$ is a modulus
of continuity of the fixed inverse,
\[
 \sup\|B_i^{-1}-B^{-1}\|
     \leq\omega_{B^{-1}}
                     \bigl(\sup\|B_i-B\|\bigr).
                                                               \tag{C.84.11}
\]
This follows by applying $B^{-1}$
to $B(B_i^{-1}z)$ and $z$.
Consequently $F\mapsto V_F$ is
continuous in the full actual
group topology. The actual group
operations make $F\mapsto H_F$
continuous as well.

On either fiber subgroup the
weighted forward difference of
the interpolations at $(r,F)$
and $(s,H)$ is bounded by
$r\|F-H\|_*+|r-s|\|H-1\|_*$.
Equations (C.84.8) and (C.84.11)
therefore give joint continuity
of each factor in the original
map-and-inverse topology.
Their product is jointly
continuous. This proves the
contraction of the entire
subspace, rather than only an
individual pathwise construction.

\emph{The explicit scope witness.}
For each integer $k$, take the
disk of radius $1/10$ centered
at $(k+1/2,0)$. These disks are
pairwise disjoint. In polar
coordinates about each center
use the identical homeomorphism
\[
 (\rho,\theta)\longmapsto(\rho,\theta+\alpha(\rho)),
 \qquad
 \alpha(\rho)=
 \begin{cases}
  \pi,&0\leq\rho\leq1/20,\\
  \pi(2-20\rho),&1/20\leq\rho\leq1/10.
 \end{cases}                                                    \tag{C.84.12}
\]
Use the identity elsewhere.
The inverse subtracts the same
angle, and the boundary values
agree. Local finiteness proves
that this is a plane homeomorphism.
It commutes with the original
unit translation and its
displacement is at most $1/5$.
Hence it belongs to $\mathscr P$.
On the vertical line through
any core center it sends the
heights $-1/40,1/40$ to
$1/40,-1/40$, respectively.
It consequently lies outside
$\mathscr M$.

The contraction (C.84.10) retains
every original operator and
boundary label on its stated
subspace. C.85--C.87 below
construct the continuous cutoff
for arbitrary original families
and a contraction on the whole
group $\mathscr P$.
\end{proof}
